% Options for packages loaded elsewhere
\PassOptionsToPackage{unicode}{hyperref}
\PassOptionsToPackage{hyphens}{url}
\PassOptionsToPackage{dvipsnames,svgnames,x11names}{xcolor}
\documentclass[
]{article}
\usepackage{xcolor}
\usepackage[a4paper, hmargin=2.2cm, vmargin=2.4cm, headheight=14.5pt, footskip=0.9cm, includehead, includefoot]{geometry}
\usepackage{amsmath,amssymb}
\setcounter{secnumdepth}{-\maxdimen} % remove section numbering
\usepackage{iftex}
\ifPDFTeX
  \usepackage[T1]{fontenc}
  \usepackage[utf8]{inputenc}
  \usepackage{textcomp} % provide euro and other symbols
\else % if luatex or xetex
  \usepackage{unicode-math} % this also loads fontspec
  \defaultfontfeatures{Scale=MatchLowercase}
  \defaultfontfeatures[\rmfamily]{Ligatures=TeX,Scale=1}
\fi
\usepackage{lmodern}
\ifPDFTeX\else
  % xetex/luatex font selection
\fi
% Use upquote if available, for straight quotes in verbatim environments
\IfFileExists{upquote.sty}{\usepackage{upquote}}{}
\IfFileExists{microtype.sty}{% use microtype if available
  \usepackage[]{microtype}
  \UseMicrotypeSet[protrusion]{basicmath} % disable protrusion for tt fonts
}{}
\makeatletter
\@ifundefined{KOMAClassName}{% if non-KOMA class
  \IfFileExists{parskip.sty}{%
    \usepackage{parskip}
  }{% else
    \setlength{\parindent}{0pt}
    \setlength{\parskip}{6pt plus 2pt minus 1pt}}
}{% if KOMA class
  \KOMAoptions{parskip=half}}
\makeatother
\usepackage{longtable,booktabs,array}
\newcounter{none} % for unnumbered tables
\usepackage{calc} % for calculating minipage widths
% Correct order of tables after \paragraph or \subparagraph
\usepackage{etoolbox}
\makeatletter
\patchcmd\longtable{\par}{\if@noskipsec\mbox{}\fi\par}{}{}
\makeatother
% Allow footnotes in longtable head/foot
\IfFileExists{footnotehyper.sty}{\usepackage{footnotehyper}}{\usepackage{footnote}}
\makesavenoteenv{longtable}
\setlength{\emergencystretch}{3em} % prevent overfull lines
\providecommand{\tightlist}{%
  \setlength{\itemsep}{0pt}\setlength{\parskip}{0pt}}
% ============================================================
% 土力学速查手册 —— 学术黑白 A4 版式（xelatex / pandoc）
% 主题：严谨、黑白、无彩色、无装饰
% ============================================================
\usepackage[fontset=windows]{ctex}            % 中文排版（宋体正文/黑体标题，自动伪粗体）
% 把常用符号划入 CJK 字体（SimSun），避免落到拉丁字体缺字：
% 带圈数字 ①-⑩、几何符号 ◆■□、箭头 →、希腊字母 αφ、数学符号 ≈
\xeCJKDeclareCharClass{CJK}{"2460->"24FF}
\xeCJKDeclareCharClass{CJK}{"25A0->"25FF}
\xeCJKDeclareCharClass{CJK}{"2190->"21FF}
\xeCJKDeclareCharClass{CJK}{"0370->"03FF}
\xeCJKDeclareCharClass{CJK}{"2248->"2248}% ≈（SimSun 有；− 等保留拉丁字体）
\usepackage{amsmath, amssymb}                 % 数学符号
\usepackage{booktabs}                         % 三线表
\usepackage{longtable}                        % 跨页表格
\usepackage{array}
\usepackage{fancyhdr}
\usepackage{enumitem}
\setlist{itemsep=2pt, topsep=4pt, leftmargin=1.6em}

% ---------- 页眉页脚：简洁学术 ----------
\pagestyle{fancy}
\fancyhf{}
\fancyhead[C]{\small\kaishu 土力学速查手册}
\fancyfoot[C]{\small\thepage}
\renewcommand{\headrulewidth}{0.4pt}
\renewcommand{\footrulewidth}{0pt}
\fancypagestyle{plain}{
  \fancyhf{}
  \fancyfoot[C]{\small\thepage}
  \renewcommand{\headrulewidth}{0pt}
}

% ---------- 表格：黑白三线风格 ----------
\renewcommand{\arraystretch}{1.35}
\setlength{\aboverulesep}{2pt}
\setlength{\belowrulesep}{2pt}

% ---------- 其他 ----------
\raggedbottom
\usepackage{bookmark}
\IfFileExists{xurl.sty}{\usepackage{xurl}}{} % add URL line breaks if available
\urlstyle{same}
% fallback for those not using the hyperref driver hyperxmp:
\makeatletter
\@ifundefined{xmpquote}{\newcommand{\xmpquote}[1]{#1}}{}
\makeatother
\hypersetup{
  colorlinks=true,
  linkcolor={black},
  filecolor={Maroon},
  citecolor={black},
  urlcolor={black},
  pdfcreator={LaTeX via pandoc}}

\author{}
\date{}

\begin{document}

\section{第六章
土压力与挡土墙}\label{ux7b2cux516dux7ae0-ux571fux538bux529bux4e0eux6321ux571fux5899}

\begin{quote}
\textbf{本章导览}：核心计算是\textbf{朗肯土压力}（主动/被动，无黏性土与黏性土两类分布），考研常结合挡土墙稳定性（抗倾覆/抗滑）出综合题。库仑土压力主要考公式与\textbf{适用条件辨析}。务必分清三种土压力大小关系与适用条件。
\end{quote}

\begin{center}\rule{0.5\linewidth}{0.5pt}\end{center}

\subsection{6.1 三种土压力}\label{ux4e09ux79cdux571fux538bux529b}

\subsubsection{公式
6-1　土压力类型与大小关系}\label{ux516cux5f0f-6-1-ux571fux538bux529bux7c7bux578bux4e0eux5927ux5c0fux5173ux7cfb}

\[E_a<E_0<E_p\]

\[E_0=K_0\gamma H\qquad(\text{分布强度 }p_0=K_0\gamma z)\]

\begin{itemize}
\tightlist
\item
  \textbf{参数说明}：\(E_a,E_0,E_p\)
  为单位墙长的主动、静止、被动土压力合力（\(\mathrm{kN/m}\)）；\(H\)
  挡墙高度；\(K_0\) 静止土压力系数；\(\gamma z\) 竖向自重应力。
\item
  \textbf{触发条件}：墙体\textbf{离开土体}------主动状态（达
  \(E_a\)）；墙体\textbf{挤压土体}------被动状态（达 \(E_p\)）。
\item
  \textbf{注意}：仅当墙有微小位移的方向与大小对应时，土压力才趋近对应值；无位移墙体用静止土压力。
\end{itemize}

\subsubsection{公式
6-2　静止土压力系数}\label{ux516cux5f0f-6-2-ux9759ux6b62ux571fux538bux529bux7cfbux6570}

\[K_0=1-\sin\varphi'\quad(\text{正常固结土，Jaky 经验式})\]

\begin{itemize}
\tightlist
\item
  \textbf{参数说明}：\(\varphi'\) 为有效内摩擦角；其他经验式常取
  \(K_0=0.5\sim0.7\)（砂土约 0.4，黏土约 0.5）。
\item
  \textbf{注意}：\(K_0\)
  为\textbf{静止条件}独有；朗肯/库仑土压力系数见下。
\end{itemize}

\begin{center}\rule{0.5\linewidth}{0.5pt}\end{center}

\subsection{6.2
朗肯土压力理论}\label{ux6717ux80afux571fux538bux529bux7406ux8bba}

\textbf{适用条件（背记）}：墙背\textbf{竖直}、墙背\textbf{光滑}（墙土摩擦角
\(\delta=0\)）、填土面\textbf{水平}。 墙背土单元处于极限平衡：主动时
\(\sigma_3\) 沿水平向，被动时 \(\sigma_1\) 沿水平向。

\subsubsection{公式
6-3　主动土压力系数与被主动系数}\label{ux516cux5f0f-6-3-ux4e3bux52a8ux571fux538bux529bux7cfbux6570ux4e0eux88abux4e3bux52a8ux7cfbux6570}

\[K_a=\tan^2\left(45^{\circ}-\frac{\varphi}{2}\right),\qquad K_p=\tan^2\left(45^{\circ}+\frac{\varphi}{2}\right)\]

\[K_a\cdot K_p=1\]

\begin{itemize}
\tightlist
\item
  \textbf{注意}：\(K_a<K_0<K_p\)；\(K_p=1/K_a\) 可快速互推。
\end{itemize}

\subsubsection{公式
6-4　无黏性土主动土压力（三角形分布）}\label{ux516cux5f0f-6-4-ux65e0ux9ecfux6027ux571fux4e3bux52a8ux571fux538bux529bux4e09ux89d2ux5f62ux5206ux5e03}

\[p_a=\gamma zK_a,\qquad E_a=\frac{1}{2}\gamma H^2K_a\]

破坏面与水平面夹角 \(45^{\circ}+\varphi/2\)；合力作用点距墙底 \(H/3\)。

\subsubsection{公式
6-5　黏性土主动土压力（含拉应力区与临界深度）}\label{ux516cux5f0f-6-5-ux9ecfux6027ux571fux4e3bux52a8ux571fux538bux529bux542bux62c9ux5e94ux529bux533aux4e0eux4e34ux754cux6df1ux5ea6}

\[p_a=\gamma zK_a-2c\sqrt{K_a}\]

\[z_0=\frac{2c}{\gamma\sqrt{K_a}}\quad(\text{临界深度，拉应力区 }0\sim z_0)\]

\[E_a=\frac{1}{2}\gamma(H-z_0)^2K_a\]

\begin{itemize}
\tightlist
\item
  \textbf{参数说明}：\(z_0\) 为墙顶拉裂区深度（m）；\(E_a\)
  合力即三角形区域（从 \(z_0\) 到 \(H\)）面积，作用点在距墙底
  \((H-z_0)/3\) 处。
\item
  \textbf{注意}：沿深度 \(0\sim z_0\)
  段土压力为负（拉应力），\textbf{不考虑它，从 \(z_0\)
  起积分}------作图与取合力时别把负段算进去。
\item
  若墙顶填土面有均布超载 \(q\)：土压力分布整体上移一个矩形
  \(\sigma_a=qK_a\)（梯形），零压点（拉应力区底缘）深度修正为：
\end{itemize}

\[z_0'=\frac{2c\sqrt{K_a}-qK_a}{\gamma K_a}=\frac{2c-q\sqrt{K_a}}{\gamma\sqrt{K_a}}\quad(q<\frac{2c}{\sqrt{K_a}}\ \text{时 }z_0'>0)\]

合力：\(E_a=\dfrac{1}{2}\gamma(H-z_0')^2K_a+qK_a(H-z_0')\)，即「三角形 +
矩形」两段相加。

\subsubsection{公式
6-6　无黏性土被动土压力}\label{ux516cux5f0f-6-6-ux65e0ux9ecfux6027ux571fux88abux52a8ux571fux538bux529b}

\[p_p=\gamma zK_p,\qquad E_p=\frac{1}{2}\gamma H^2K_p\]

合力作用点距墙底 \(H/3\)，三角形分布。

\subsubsection{公式
6-7　黏性土被动土压力}\label{ux516cux5f0f-6-7-ux9ecfux6027ux571fux88abux52a8ux571fux538bux529b}

\[p_p=\gamma zK_p+2c\sqrt{K_p}\]

\[E_p=\frac{1}{2}\gamma H^2K_p+2cH\sqrt{K_p}\]

梯形分布（由 \(z=0\) 处 \(2c\sqrt{K_p}\)
起），不能用简单的「三角形中心」取合力位置------要按梯形求形心。

\subsubsection{公式
6-8　朗肯土压力计算步骤汇总（计算题模板）}\label{ux516cux5f0f-6-8-ux6717ux80afux571fux538bux529bux8ba1ux7b97ux6b65ux9aa4ux6c47ux603bux8ba1ux7b97ux9898ux6a21ux677f}

\begin{enumerate}
\def\labelenumi{\arabic{enumi}.}
\tightlist
\item
  判断墙型条件是否满足朗肯（竖直、光滑、填土面水平）；
\item
  查/算 \(K_a,K_p\)；
\item
  区分无黏性土（\(c=0\)，三角形）或黏性土（先算 \(z_0\)，三角形/梯形）；
\item
  分布图 → 面积（合力 \(E\)），标注作用点位置；
\item
  若有地下水：水位以上用 \(\gamma\)、以下用 \(\gamma'\)
  且再加\textbf{静水压力}
  \(p_w=\gamma_w z_w\)（水压力单独叠加，作用于墙身）。
\end{enumerate}

\begin{quote}
\textbf{注意}：有地下水的挡土墙，土压力与水压力必须\textbf{分算再叠加}，不能用
\(\gamma_{sat}\) 直接乘------这是考研综合题的常扣分点。
\end{quote}

\begin{center}\rule{0.5\linewidth}{0.5pt}\end{center}

\subsection{6.3
库仑土压力理论}\label{ux5e93ux4ed1ux571fux538bux529bux7406ux8bba}

\textbf{适用条件（背记）}：墙背可\textbf{倾斜}、可\textbf{粗糙（\(\delta\ne0\)）}、填土面可\textbf{倾斜
\(\beta\)}；适用无黏性土（也可推广）。核心思路是\textbf{滑动楔体极限平衡}。

\subsubsection{公式
6-9　库仑主动土压力系数（一般式）}\label{ux516cux5f0f-6-9-ux5e93ux4ed1ux4e3bux52a8ux571fux538bux529bux7cfbux6570ux4e00ux822cux5f0f}

\[K_a=\frac{\sin^2(\alpha+\varphi)}{\sin^2\alpha\,\sin(\alpha-\delta)\left(1+\sqrt{\dfrac{\sin(\varphi+\delta)\sin(\varphi-\beta)}{\sin(\alpha-\delta)\sin(\alpha+\beta)}}\right)^2}\]

\[E_a=\frac{1}{2}\gamma H^2K_a\]

\begin{itemize}
\tightlist
\item
  \textbf{参数说明}：\(\alpha\) 为墙背与水平面夹角（竖直墙
  \(\alpha=90^{\circ}\)）；\(\delta\) 为墙背与填土间摩擦角；\(\beta\)
  为填土面倾角；\(\varphi\) 为填土内摩擦角。
\item
  \textbf{退化验证}：令
  \(\alpha=90^{\circ},\ \delta=0,\ \beta=0\)，上式化为
  \(K_a=\tan^2(45^{\circ}-\varphi/2)\)------朗肯是库仑的特例。
\item
  \textbf{作用点}：无黏性土（三角形分布）作用于距墙底
  \(H/3\)；合力方向与墙背法线成 \(\delta\) 角（不水平！）。
\item
  \textbf{注意}：计算中角度取\textbf{弧度}或角度制要统一；\(\delta\)
  的符号工程上常偏安全取负。
\end{itemize}

\subsubsection{公式
6-10　库仑被动土压力}\label{ux516cux5f0f-6-10-ux5e93ux4ed1ux88abux52a8ux571fux538bux529b}

\[K_p=\frac{\sin^2(\alpha-\varphi)}{\sin^2\alpha\,\sin(\alpha+\delta)\left(1-\sqrt{\dfrac{\sin(\varphi+\delta)\sin(\varphi+\beta)}{\sin(\alpha+\delta)\sin(\alpha+\beta)}}\right)^2},\quad E_p=\frac{1}{2}\gamma H^2K_p\]

\begin{itemize}
\tightlist
\item
  \textbf{注意}：\(K_p\) 的 \(\varphi\)
  全变号、分母括号取减号------公式形似而符号相反，考场照抄易错，建议记住括号里是
  \(\varphi+\beta\) 且根号前置负号。
\end{itemize}

\begin{center}\rule{0.5\linewidth}{0.5pt}\end{center}

\subsection{6.4
挡土墙的稳定性验算}\label{ux6321ux571fux5899ux7684ux7a33ux5b9aux6027ux9a8cux7b97}

设墙每米自重为 \(W\)（作用线距墙趾 \(x_0\)），主动土压力合力
\(E_a\)（水平向分量），\(E_a\) 作用点距墙趾竖直高度 \(y_0\)。

\subsubsection{公式
6-11　抗倾覆稳定}\label{ux516cux5f0f-6-11-ux6297ux503eux8986ux7a33ux5b9a}

\[K_t=\frac{W\,x_0}{E_a\,y_0}\ge1.5\]

\begin{itemize}
\tightlist
\item
  \textbf{参数说明}：规范要求 \(K_t\ge1.5\)（部分教材取
  1.6）；分母为土压力使墙绕墙趾 A
  的倾覆力矩。若挡土墙为混凝土，墙背土压力的水平分量为
  \(E_{ax}=E_a\cos(\delta-\alpha')\)，竖直分量为
  \(E_{ay}=E_a\sin(\delta-\alpha')\)，用水平分量求倾覆矩即可。
\item
  \textbf{注意}：\(x_0,y_0\)
  都取到墙趾（转动点）的力臂；这是最常配错误的正负号题。
\end{itemize}

\subsubsection{公式
6-12　抗滑稳定}\label{ux516cux5f0f-6-12-ux6297ux6ed1ux7a33ux5b9a}

\[K_s=\frac{(W+E_{ay})\mu}{E_{ax}}\ge1.3\]

\begin{itemize}
\tightlist
\item
  \textbf{参数说明}：\(\mu\) 为墙底与地基摩擦系数；\(E_{ax},E_{ay}\)
  为主动土压力水平、竖向分量；规范要求 \(K_s\ge1.3\)。
\item
  \textbf{注意}：竖向分量 \(E_{ay}\)
  增大了墙底正压力、有助于抗滑，别漏；水平分量才是滑动力。
\end{itemize}

\subsubsection{公式
6-13　挡墙基底压力验算}\label{ux516cux5f0f-6-13-ux6321ux5899ux57faux5e95ux538bux529bux9a8cux7b97}

\[p_{max}=\frac{N}{A}\left(1+\frac{6e}{B}\right)\le1.2f\]

\[p_{min}=\frac{N}{A}\left(1-\frac{6e}{B}\right)\ge0,\qquad e=\frac{B}{2}-\frac{M_{\text{稳}}-M_{\text{倾}}}{N}\]

\begin{itemize}
\tightlist
\item
  \textbf{参数说明}：\(N\) 为基底竖向合力；\(e\) 偏心距；\(B\)
  基底宽；\(f\) 持力层承载力特征值。判定同第三章公式 3-6/3-7。
\end{itemize}

\begin{center}\rule{0.5\linewidth}{0.5pt}\end{center}

\subsection{6.5
常见挡土墙类型（概念题）}\label{ux5e38ux89c1ux6321ux571fux5899ux7c7bux578bux6982ux5ff5ux9898}

\begin{itemize}
\tightlist
\item
  \textbf{重力式}：依靠自重抗滑抗倾覆，造价低，适用墙高不大；回填要求低，应用最广。
\item
  \textbf{悬臂式}：钢筋混凝土，利用底板土重，适用中等墙高。
\item
  \textbf{扶壁式}：悬臂墙+扶壁，适用高墙。
\item
  结构形式选择与「挡墙高 + 地基土质」相关------多为判断题/简答题。
\end{itemize}

\begin{center}\rule{0.5\linewidth}{0.5pt}\end{center}

\subsection{6.6
本章公式索引小结}\label{ux672cux7ae0ux516cux5f0fux7d22ux5f15ux5c0fux7ed3}

{\def\LTcaptype{none} % do not increment counter
\begin{longtable}[]{@{}
  >{\raggedright\arraybackslash}p{(\linewidth - 4\tabcolsep) * \real{0.3333}}
  >{\raggedright\arraybackslash}p{(\linewidth - 4\tabcolsep) * \real{0.3333}}
  >{\raggedright\arraybackslash}p{(\linewidth - 4\tabcolsep) * \real{0.3333}}@{}}
\toprule\noalign{}
\begin{minipage}[b]{\linewidth}\raggedright
编号
\end{minipage} & \begin{minipage}[b]{\linewidth}\raggedright
公式
\end{minipage} & \begin{minipage}[b]{\linewidth}\raggedright
用途
\end{minipage} \\
\midrule\noalign{}
\endhead
\bottomrule\noalign{}
\endlastfoot
6-1 & \(E_a<E_0<E_p\) & 三种土压力关系 \\
6-2 & \(K_0=1-\sin\varphi'\) & 静止土压力系数 \\
6-3 &
\(K_a=\tan^2(45^{\circ}-\varphi/2),\ K_p=\tan^2(45^{\circ}+\varphi/2)\)
& 朗肯系数 \\
6-4 & \(p_a=\gamma zK_a,\ E_a=\dfrac{1}{2}\gamma H^2K_a\) &
无黏性土主动 \\
6-5 &
\(p_a=\gamma zK_a-2c\sqrt{K_a},\ z_0=\dfrac{2c}{\gamma\sqrt{K_a}}\) &
黏性土主动 \\
6-6 & \(E_p=\dfrac{1}{2}\gamma H^2K_p\) & 无黏性土被动 \\
6-7 & \(E_p=\dfrac{1}{2}\gamma H^2K_p+2cH\sqrt{K_p}\) & 黏性土被动 \\
6-9 & 库仑 \(K_a\) 一般式 & 倾斜/粗糙墙背 \\
6-10 & 库仑 \(K_p\) 一般式 & 被动状态（符号易混） \\
6-11 & \(K_t=\dfrac{Wx_0}{E_ay_0}\ge1.5\) & 抗倾覆 \\
6-12 & \(K_s=\dfrac{(W+E_{ay})\mu}{E_{ax}}\ge1.3\) & 抗滑 \\
6-13 & \(p_{max}=\dfrac{N}{A}(1+\dfrac{6e}{B})\le1.2f\) &
挡墙基底压力验算 \\
\end{longtable}
}

\begin{quote}
\textbf{考试高频提醒}：① 朗肯适用条件脱口而出：竖直、光滑、填土面水平；②
黏性土主动取合力时\textbf{从 \(z_0\) 开始}、别算拉力区；③
有水时土压力与水压力\textbf{分算叠加}；④ 库仑 \(K_a\) 公式在
\(\alpha=90^{\circ},\delta=0,\beta=0\) 时退化为朗肯，可自检；⑤
抗倾覆对\textbf{墙趾}取矩，抗滑用\textbf{水平分量}。
\end{quote}

\end{document}
